\section{Diseño Visual}
\label{sec:visualdesign}

Con base en los objetivos de diseño (G1--G4, Sección~\ref{sec:design-goals}) y el análisis de tareas (T1--T8, Sección~\ref{sec:task-analysis}), se propone un sistema de tres vistas coordinadas, cada una compuesta por sub-paneles vinculados. La estructura de esta sección --Descripción seguida de Justificación por vista, y una subsección final de Interacciones-- sigue el esquema de presentación adoptado en trabajos previos de análisis visual multimodal~\cite{Emoco2019}.

\subsection{Vista A: Espacio de Representaciones Fusionadas}
\label{sec:vista_a}

\textbf{Descripción.} Esta vista opera sobre las representaciones fusionadas de Husformer (\texttt{last\_hs}) agregadas a nivel de trial. Está compuesta por tres sub-paneles: (A1) una proyección de baja dimensionalidad (PCA, UMAP o t-SNE, seleccionable por el usuario) del espacio de representación; (A2) una capa de agrupamiento algorítmico (K-Means o HDBSCAN) superpuesta sobre la proyección, codificada por color; y (A3) un panel de comparación de perfil, vinculado a la selección múltiple
realizada en A1: al marcar uno o varios trials de interés, A3 agrupa
automáticamente sus participantes correspondientes (sin duplicados -- varios
trials de un mismo participante producen una única fila) y despliega, para
ese subconjunto, sus atributos de cuestionario (categóricos y numéricos,
p.\ ej.\ género, lateralidad manual, consumo de alcohol/cafeína, edad, horas
de sueño) mediante una codificación de barras compactas por atributo,
inspirada en LineUp~\cite{gratzl2013lineup}.

\textbf{Justificación.} A1 atiende directamente a T1 (identificar agrupamientos o valores atípicos en el espacio de representación), replicando y extendiendo el enfoque exploratorio ya validado en el análisis de características manuales (Tarea 1). A2 se incorpora para atender T2 (comparar trials/participantes de forma estructurada): si bien la proximidad visual en una proyección ya sugiere similitud, una capa de agrupamiento algorítmico explícito reduce la ambigüedad perceptual de "leer" clusters a simple vista, un enfoque con precedente en sistemas de análisis visual interactivo que combinan proyecciones con clustering algorítmico para apoyar tareas de comparación~\cite{evoair}. A3 no compara los trials/participantes seleccionados dentro del espacio de
representación en sí --esa comparación estructural la habilita A2 mediante
el agrupamiento algorítmico-- sino que ofrece una vía complementaria a T1:
una vez identificado un trial o participante cuya representación se aparta
del resto, A3 permite examinar si ese subconjunto comparte rasgos
demográficos o de cuestionario, aportando una hipótesis explicativa
adicional más allá del espacio de representación en sí. La codificación
visual de A3 sigue el patrón de tablas de comparación multi-atributo de
LineUp~\cite{gratzl2013lineup} (barras compactas por celda en lugar de texto
plano).

\subsection{Vista B: Atención Temporal del Trial}
\label{sec:vista_b}

\textbf{Descripción.} Al seleccionar un trial (desde A1--A3), esta vista despliega su dinámica interna de atención a lo largo del tiempo, mediante dos sub-paneles: (B1) un mapa de calor de dominancia de modalidad sobre las ventanas temporales del trial; y (B2) un panel de comparación de señales fisiológicas crudas normalizadas, seleccionables por modalidad (EEG por región anatómica o hemisferio; EOG/EMG/GSR/Resp+Plet+Temp por canal), superpuestas a lo largo de todo el trial.

\textbf{Justificación.} B1 atiende conjuntamente a T3 (explorar la dinámica temporal de atención dentro de un trial) y T4 (identificar segmentos temporales donde una modalidad domina la representación fusionada): la estructura de matriz ofrece de un vistazo tanto la evolución temporal completa (overview, T3) como la identificación puntual de qué modalidad domina en cada ventana mediante intensidad de color (T4), siguiendo la máxima "overview primero, zoom y filtro, detalles bajo demanda"~\cite{shneiderman1996}. B2 atiende a T5 (relacionar picos o cambios abruptos en la atención con eventos visibles en la señal original), permitiendo comparar señales fisiológicas crudas de distintas modalidades o regiones, normalizadas y superpuestas, a lo largo de todo el trial. Al estar siempre visible junto a B1 dentro de la misma Vista~B, el usuario puede relacionar un pico de atención observado en B1 con un evento correspondiente en la señal cruda de B2 sin depender de la memoria, sin necesidad de duplicar la representación de la atención dentro del propio B2 -- un resaltado sincronizado bidireccional entre ambos paneles (\textit{linked highlighting}~\cite{becker1987brushing}) refuerza esta relación sin reconstruir ninguno de los dos.

\textbf{Nota de proceso.} Una versión anterior de esta vista incluía un tercer sub-panel entre el heatmap y la comparación de señal cruda -- una serie temporal agregada por modalidad mediante líneas, mostrando el mismo dato que B1 (dominancia por modalidad) con un idioma visual distinto. Se descartó por decisión de diseño: no atendía ninguna tarea que B1 no cubriera ya, y liberar su espacio permitió ampliar el panel de señal cruda (B2), que se beneficia de más ancho para su selector de canales y su zoom temporal.

\subsection{Vista C: Detalle Anclado a la Señal Cruda}
\label{sec:vista_c}

\textbf{Descripción.} A diferencia del resto del sistema, Vista~C no depende de una selección hecha en Vista~A ni de un trial cargado en Vista~B por sí solo: se dispara al pasar el cursor sobre la señal cruda de \textbf{B2}, mostrando el detalle correspondiente a ese instante puntual mediante dos sub-paneles: (C1) la matriz de atención cross-modal $5\times5$ (\texttt{attn\_cross\_summary}) de la ventana bajo el cursor, sin promediar; y (C2), para cada modalidad activa en B2, dos mini-gráficos apilados (no superpuestos en un eje compartido) sobre una ventana de $\pm3$ segundos alrededor del cursor: la señal real sin normalizar, y el porcentaje de dominancia de esa modalidad (mismo dato que B1). Una línea guía vertical sincroniza la posición exacta entre ambos gráficos de cada tarjeta.

\textbf{Justificación.} C1 atiende a T6 (inspeccionar qué modalidad recibe mayor peso de atención cross-modal en un instante dado), ahora anclado a un instante elegido observando directamente la señal fisiológica real, en vez de una ventana elegida sobre un dato ya derivado (el heatmap de B1) -- una conexión más directa con el objetivo real de la tarea. C2 atiende a T8 (comparar la señal cruda de una modalidad frente a su peso de atención recibido en un instante dado), validando si la atención del modelo coincide con el conocimiento de dominio esperado sobre la relación señal-emoción; se optó por gráficos apilados y no por un eje compartido (\textit{superimpose}), para no sugerir visualmente una relación de escala inexistente entre unidades físicas distintas ($\mu V$/$\mu S$ de la señal frente a porcentaje de dominancia).

\textbf{Nota de proceso.} Esta vista pasó por tres diseños durante el desarrollo: el original (descrito en versiones previas de este documento), un drill-down desde una ventana seleccionada en B1/B2 con un tercer sub-panel dedicado a comparar varias ventanas entre sí (ver nota sobre T7 retirada, \cref{sec:task-analysis}); un diseño intermedio de Small Multiples sobre los trials seleccionados en Vista~A, implementado por completo y descartado por decisión de diseño; y el vigente, descrito arriba, anclado a una acción sobre Vista~B en vez de una selección de Vista~A o un click sobre Vista~B.

\subsection{Interacciones}
\label{sec:interacciones}

Las tres vistas están coordinadas mediante los siguientes mecanismos de interacción, en línea con el principio de \textit{multiple linked views} ampliamente adoptado en sistemas de análisis visual multimodal~\cite{Emoco2019}:

\begin{itemize}
    \item \textbf{Clicking:} seleccionar un punto en A1/A2 carga el trial correspondiente en Vista~B (B1 y B2).
    \item \textbf{Selección múltiple:} click en A1/A2 agrega o quita un trial de una selección compartida entre ambos paneles; A3 se actualiza automáticamente con los participantes correspondientes a esa selección.
    \item \textbf{Hovering sincronizado (\textit{linked highlighting}, B1$\leftrightarrow$B2):} pasar el cursor sobre una ventana puntual en B1 o en B2 resalta la misma ventana en el otro panel, sin reconstruir ninguno de los dos~\cite{becker1987brushing}.
    \item \textbf{Hovering (drill-down B2$\to$C1/C2):} pasar el cursor sobre la señal cruda en B2 actualiza en tiempo real tanto la matriz mostrada en C1 como el detalle de C2, permitiendo recorrer varios instantes consecutivos sin una acción discreta por cada uno -- relevante porque las matrices de atención cross-modal cambian de forma sutil entre ventanas vecinas. C1 y C2 conservan el último instante mostrado al retirar el cursor (no vuelven a un estado vacío), de modo que la vista permanece disponible mientras el usuario la examina.
    \item \textbf{Linking:} las tres vistas conforman un flujo de exploración progresiva (\textit{drill-down}) de participante/trial (A) a dinámica del trial (B) a detalle anclado a la señal cruda (C).
    \item \textbf{Filtrado global:} filtros por participante y por trial en A1, persistentes mientras no se cambien explícitamente.
    \item \textbf{Details-on-demand:} información numérica exacta disponible mediante \textit{tooltips} en A1, B1, C1 y C2.
\end{itemize}